
%%%%%%%%%%%%%%%%%%%%%%%%%%%%%%%%%%%%%%%%%%%%%%%%%%%%%%%%%%%%%%%%%%%%%%%%%%%%%%%%%%%%%%%%%%%%%%%%%%% 
% KAPITOLA 8
%%%%%%%%%%%%%%%%%%%%%%%%%%%%%%%%%%%%%%%%%%%%%%%%%%%%%%%%%%%%%%%%%%%%%%%%%%%%%%%%%%%%%%%%%%%%%%%%%%

\chapter{Vyhodnotenie}
\label{ch:kap8}
Táto kapitola sa zameriava na vyhodnotenie výsledkov experimentov, ktoré boli definované v Kapitole~\ref{ch:kap7} na otázky RQ1-RQ6. Ako prvé sa zameriava na vyhodnotenie výsledkov jadra testovania, modelu qwen2.5 v troch konfiguráciách. Testovanie sa zameriava na testovanie modelu vo veľkostiach 3b, 7b a zmiešanú konfiguráciu oboch veľkostí, naprieč šiestimi teplotami. Ďalej sa zameriava na výsledky rozšírenia kde cieľom je porovnanie modelových rodín. Medzi modelové rodiny patrí llama3.1:8b a mistral:7b s referenčnou konfiguráciou qwen2.5. Kapitola sa zameriava na analýzu zlyhaní a zhrnutie limitov vyhodnotenia. Na záver kapitoly je demonštrácia, teda ukážka a overenie, že nástroj je schopný fungovať aj nad inými DBMS ako MySQL. DBMS vybrané na porovnanie je PostgreSQL. 

Všetky tri úrovne sa zameriavajú na meranie technickej použiteľnosti receptu. Sémantickú správnosť  hodnôt nemerajú. Riadok môže úspešne prejsť L3 aj v tom prípade, že obsahuje hodnotu, ktorá pre daný stĺpec nedáva zmysel.


\section{Priebeh experimentov}
Vyhodnotenie prebieha podľa testovacích scenárov v troch fázach: jadro (31 databáz), rozšírenie (19 základných databáz) a demonštrácia (10 databáz). Každá konfigurácia bola spustená v troch kolách, pričom opakované behy ukázali, že ani teplota 0.0 nezaručuje rovnaké výstupy. Práve preto sú všetky behy vyhodnocované v troch kolách a výsledky sa ukazujú spoločne za všetky kolá.


Ako dodatok sa uvádza úspešnosť L3\_Adjusted ako v dvoch agregáciách. Jedna ako celková úspešnosť, teda podiel úspešných príkazov zo všetkých vygenerovaných a druhá ako priemerná úspešnosť po databázach. Dôvodom je, že väčšie databázy prispievajú k celkovej úspešnosti výrazne viac a preto je tam zavedená aj druhá miera úspešnosti, kde má každá databáza rovnakú váhu. Takto je aj možné vidieť ako sú rozložené zlyhania databáz, t.j., či sú rozložené rovnomerne alebo sa chyby ukazujú len v niektorých databázach. Keďže úspešnosti sa pohybujú v extrémnych rozdieloch 0\% a 100\%, tak sa používa Wilcoxonov test na porovnanie tej istej databázy pri dvoch teplotách.


Počas behu testov došlo k výpadku dvoch behov konfigurácie qwen2.5:7b (testdb18 pri teplote 0.5 a testdb9 pri teplote 1.2). Výpadok je zaznamenaný v tabuľke~\ref{tab:vysledky_jadro}.

\subsection{Prínos opravnej vrstvy}
Je taktiež dôležité spomenúť, že validačná a opravná vrstva (Sekcia~\ref{sub:validacia_oprava}) nie je vo výsledkoch explicitne spomenutá. Pre ukážku boli zhodnotené opravy na referenčnej konfigurácii (zmiešaná konfigurácia, qwen2.5, teplota 0.0, MySQL) nad 31 databázami v troch kolách (Tabuľka~\ref{tab:opravy}).


\begin{table}[H]
    \centering
    \begin{tabular}{l r}
    \toprule
       Počet behov  & 93 \\
       Počet zásahov opravnej vrstvy & 3690 \\
       Podiel behov s aspoň jedným zásahom  & 67.7\% \\
    \bottomrule
    \end{tabular}
    \caption{Zásahy opravnej vrstvy na referenčnej konfigurácii}
    \label{tab:opravy}
\end{table}

Jednotkou merania je zásah, nie zachránený recept. Jeden recept môže vyžadovať desiatky zásahov. Čo sa týka typu zásahov, tie sú takmer výhradne hodnotové, čo ukazuje na nízky počet zlyhaní na úrovni L1. Prínos opravnej vrstvy teda nie je v záchrane štruktúry dokumentu, ale v úplnosti mapovania.

\section{Vplyv teploty (RQ1)}

\begin{table}[H]
    \centering
    \small
    \begin{tabular}{cccccc}
    \toprule
    & & & & \multicolumn{2}{c}{\textbf{L3\_Adjusted}} \\
    \textbf{Konfig.} & \textbf{teplota} & \textbf{L2} & \textbf{\%} & \textbf{celková \%} & \textbf{priemer db $\pm$ sd \%} \\
    \midrule
        3b & 0.0 & 55/93 & 59.1 & 85.4 & 83.1 $\pm$ 32.8 \\ 
        3b & 0.2 & 58/93 & 62.4 &  85.3 & 77.1 $\pm$ 38.6 \\ 
        3b & 0.3 & 57/93 & 61.3 &  90.9 & 75.9 $\pm$ 41.5 \\ 
        3b & 0.5 & 55/93 & 59.1 &  90.2 & 84.9 $\pm$ 31.9 \\ 
        3b & 0.8 & 55/93 & 59.1 & 94.4 & 87.9 $\pm$ 29.8 \\
        3b & 1.2 & 56/93 & 60.2 & 85.0 & 71.2 $\pm$ 42.1 \\
        \midrule
        7b & 0.0 & 60/93 & 64.5 &  100.0 & 100.0 $\pm$ 0.0 \\ 
        7b & 0.2 & 68/93 & 73.1 & 100.0 & 100.0 $\pm$ 0.0 \\ 
        7b & 0.3 & 66/93 & 71.0 &  100.0 & 100.0 $\pm$ 0.0 \\ 
        7b & 0.5 & 63/92 & 68.5 &  100.0 & 100.0 $\pm$ 0.0 \\ 
        7b & 0.8 & 67/93 & 72.0 &  98.8 & 97.2 $\pm$ 15.5 \\ 
        7b & 1.2 & 61/92 & 66.3 &  99.1 & 96.4 $\pm$ 18.7 \\ 
        \midrule
        mix & 0.0 & 75/93 & 80.6 &  98.1 & 95.7 $\pm$ 20.5 \\ 
        mix & 0.2 & 74/93 & 79.6 &  98.4 & 96.8 $\pm$ 17.1 \\
        mix & 0.3 & 74/93 & 79.6 &  97.6 & 94.1 $\pm$ 23.7 \\ 
        mix & 0.5 & 72/93 & 77.4 &  99.1 & 97.7 $\pm$ 13.3 \\ 
        mix & 0.8 & 68/93 & 73.1 &  99.3 & 97.8 $\pm$ 13.4 \\ 
        mix & 1.2 & 77/93 & 82.8 &  97.6 & 94.1 $\pm$ 23.7 \\ 
    \bottomrule
    \end{tabular}
    \caption{Úspešnosť konfigurácií qwen2.5 (6 teplôt, 31 databáz, 3 behy), L2 = akceptované recepty}
    \label{tab:vysledky_jadro}
\end{table}



Pri tejto otázke bolo cieľom skúmať ako veľmi má teplota vplyv na úspech v jednotlivých úrovniach, čiže správnosť YAML formátu (L1), prijatie receptu Datasterom (L2)  a spustenie SQL príkazov v databáze (L3). Testovanie prebiehalo s modelom qwen2.5 s veľkosťami modelov 7b, 3b a ich zmiešanou konfiguráciou (mix).


Zlyhania na úrovni L1 sa vyskytli len v malom počte behov (8 z 1672), čo poukazuje na to, že využitie validačnej a opravnej vrstvy programu robí štruktúrnu správnosť YAML nezávislou od teploty.


Výsledky pre vrstvu L2 sa pohybujú v priemere okolo 60\% pre qwen2.5:3b, 69\% pre qwen2.5:7b a 78\% pre mix. Vo výsledkoch nie je vidieť žiadny vplyv teploty, potvrdil to aj Wilcoxonov test na priemeroch jednotlivých databáz. Rozdiely medzi dvojicami teplôt neboli významné.  Zaujímavé ale je ako na úrovni $L3_{Adjusted}$ má teplota (1.2) najnižšie priemery a zároveň najväčšie výkyvy. 


Pre ďalšie fázy bola zvolená teplota 0.0, nie ako najlepší výsledok, ale predvídateľný štandardný základ.

\section{Vplyv veľkosti modelu (RQ2)}
\label{sub:odpoved_rq2}
Na rozdiel od teploty sa veľkosť modelu ukázalo ako veľký faktor ovplyvňujúci úspešnosť behov. Využitie validačnej a opravnej vrstvy v L1 dosiahlo vysokú priemernú úspešnosť 99.5\%. Spomedzi všetkých testov model 3b dosiahol najnižšie výsledky na úrovniach L2 a L3 (L2 priemer 60\%, L3 90\%). Pri konfigurácii 7b a mix je dosiahnutá úspešnosť okolo 98 - 100\% pre L3\_Adjusted, teda ak bol recept akceptovaný v L2 tak boli SQL dotazy vykonané takmer bez chyby.


Čo sa týka úspešnosti L2 je vidieť ako mix dosahuje vyššie čísla ako 7b (priemer 78\% oproti 69\%), toto je ale vysvetliteľné, model 7b častejšie generuje sebareferenčné odkazy FROM (stĺpec odkazuje sám na seba), ktoré validačná vrstva nezachytáva a ktoré vedú k pádu generátora. Pri modeli 7b spadlo až 49 behov, zatiaľ čo pri mix 7b a 3b len 7 a nakoniec u 3b spadlo len 5 behov. Ak tieto zlyhania odpočítame tak sa úspešnosť 7b zvýši na okolo 78\% a mix na 80\%. Záverom nie je ale ukázať, že mix  je lepší ale poznamenať ako mix dosahuje rovnakú kvalitu ako 7b napriek nižším nárokom, keďže dve zo štyroch fáz generovania obsluhuje menší model.


Najviac sa rozdiely veľkostí ale prejavujú pri prechode zo základnej a reprodukčnej sady databáz na sadu extrémnych databáz, úspešnosť modelu 3b v tomto porovnaní prudko klesá zo 68\% na 40\%, zatiaľ čo model 7b zostáva nemenný na 69\% a mix klesá iba mierne zo 82\% na 72\%. Dôležitý je aj počet nezvládnutých databáz\footnote{Databázy s úspešnosťou $\leq$ 50\%}, teda až 6 pre 3b, pri konfigurácii 7b a mix bol tento počet len 3. Zároveň bola pri modeli 3b priemerná úspešnosť pri 100 tabuľkách vyššia ako pri 60 tabuľkách, zatiaľ čo pri konfiguráciach 7b a mix zostávajú tieto pomery úplne identické.


\section{Vplyv modelových rodín (RQ3)}

Porovnanie rodín prebehlo s referenčnou teplotou 0.0 na 19 základných databázach (Tabuľka~\ref{tab:vysledky_rodina}). Čo sa týka modelu llama3.1:8b tak je možné vidieť podobné výsledky s qwen2.5:7b (L2 67\% oproti 63\% a L3\_Adjusted 95\% oproti 100\%), pričom k väčšine zlyhaní došlo kvôli dataster\_smt\_count (Tabuľka~\ref{tab:klasifikacia_zlyhani}). Týmto je zjavné, že modely porovnateľnej veľkosti  z rôznych rodín poskytujú podobnú kvalitu.


Na opačnej strane je ale možné vidieť ako mistral:7b výrazne zaostáva na oboch úrovniach (L2 42\% a L3\_Adjusted iba 29\%). Najčastejšie kategórie zlyhaní sú data\_too\_long a fk\_violation. Teda je jasné, že mistral produkuje výrazne horšie dáta v akceptovaných receptoch ako aj to, že celkový počet akceptovaných receptov je menší. 

\begin{table}[H]
    \centering
    \small
    \begin{tabular}{l r r r r} 
    \toprule
    & & &  \multicolumn{2}{c}{\textbf{L3\_Adjusted}} \\
    \textbf{Model} & \textbf{L2} & \textbf{\%}  & \textbf{celková \%} & \textbf{priemer db $\pm$ sd \%} \\\midrule
    
        qwen2.5:3b & 39/57 & 68.4 & 79.8 & 82.6 $\pm$ 31.5 \\
        qwen2.5:7b & 36/57 & 63.2 & 100.0 & 100.0 $\pm$ 0.0 \\
        qwen mix & 51/57 & 89.5 & 94.1 & 93.8 $\pm$ 24.5 \\
        llama3.1:8b & 38/57 & 66.7 & 95.1 & 94.3 $\pm$ 23.6 \\
        mistral:7b & 24/57 & 42.1 & 28.7 & 40.0 $\pm$ 41.8 \\
    \bottomrule
    \end{tabular}
    \caption{Porovnanie rodín (teplota: 0.0, 19 zákl. databáz, 3 behy)}
    \label{tab:vysledky_rodina}
\end{table}



Samotná veľkosť modelu kvalitu nezaručuje, ale rozhoduje aj rodina.

\section{Zmena úspešnosti s rastúcou zložitosťou databázovej schémy (RQ4, RQ5)}
Podiel behov, v ktorých bol recept akceptovaný nástrojom Dataster (6 teplôt $\times$ 3 behy, Tabuľka~\ref{tab:zlozitost_schemy}).

\begin{table}[H]
    \centering
    \begin{tabular}{ccccc}
        \toprule
        Typ & úroveň & 3b & 7b  & mix\\
        \midrule
         počet tabuliek & 30 & 7/18  & 18/18  &  18/18 \\ 
                       & 60 & 0/18 & 2/18 & 4/18 \\ 
                       & 100 & 9/18  & 2/18  & 4/18\\ 
        \midrule
         šírka tabuľky & 50  & 18/18  & 18/18 & 18/18\\ 
                       & 100 & 11/18  & 18/18 & 18/18\\ 
                       & 200 & 1/18   & 1/18  & 1/18\\ 
        \midrule
         hĺbka FK      & 10 & 1/18  & 18/18 &  18/18\\ 
                       & 20 & 3/18  & 17/18 &  18/18\\
        \midrule
         fan-out       & 21 & 14/18 & 18/18 & 17/18\\ \bottomrule
        \end{tabular} 
        \caption{Zmena úspešnosti pri komplexných tabuľkách}
        \label{tab:zlozitost_schemy}
\end{table}

Každý sledovaný parameter (počet tabuliek, šírka tabuliek a hĺbka FK) sa správa inak. Pri počte tabuliek klesá úspešnosť niekde medzi 30 a 60 tabuľkami. Pri 60 tabuľkách konfigurácie 7b zvládli vygenerovať recept len pre 2 behy a mix len pre 4. Zaujímavé je ale, že priemerná úspešnosť pri 100 tabuľkách je mierne vyššia ako pri 60, takže o zlyhaní nerozhoduje len počet tabuliek ale aj ich konkrétny charakter.


Pri šírke tabuliek (počte stĺpcov) funguje generovanie spoľahlivo do 100 stĺpcov a pri 200 stĺpcoch padá na úspešnosť 1/18 pre všetky konfigurácie. Dominantným zlyhaním je ale RecursionError vo vnútri nástroja Dataster, takže tento strop patrí Datasteru, nie modelom. 


Na záver čo sa týka hĺbky reťazca. Problém nerobí konfigurácia 7b a mix a nízke čísla pre 3b spôsobuje prevažne dataster\_smt\_count, čiže interakcia receptov menšieho modelu so SMT riešiteľom, nie hĺbka reťazca samotná.


Odpoveď na RQ5 sa líši v tom, že pri počte tabuliek nie je možné presnejšie určiť kde leží zlom medzi 30 a 60 tabuľkami. Pri šírke tabuliek je to medzi 100 a 200 stĺpcami a pri hĺbke FK reťazca v meranom rozsahu 20 tabuľkových úrovní sa zlom pre konfigurácie 7b a mix neukázal.

\section{Odlíšenie zlyhaní spôsobených navrhnutou pipeline od zlyhaní generátora Dataster (RQ6)}
\label{sub:k8_odpoved_rq6}
Áno, jedným z kľúčových cieľov testovania bolo spoľahlivo rozoznať priradenie chýb medzi samotným programom a generátorom Dataster. Toto sa opiera o päť nezávislých pozorovaní.


Najčastejšie sa objavoval defekt dataster\_sql\_escaping, kde bolo odhalené, že Dataster v príkazoch INSERT vygeneroval neescapované hodnoty. To, že ide o chybu generátora, potvrdil experiment, keď sa zaplo escapovanie v SQL  a tým aj zlyhania zmizli. 


Ďalej bol pozorovaný defekt dataster\_finterval\_fraction, kde pre hodnoty finterval boli dosadené zlomky namiesto desatinných čísel. Toto hovorí o defekte generátora  a potvrdzuje to aj podobné správanie dvoch DBMS (MySQL a PostgreSQL), tieto hodnoty boli odmietnuté, chyba teda vznikla pred databázou.


Je dôležité spomenúť aj kategóriu smt\_count. Priradenie zlyhania k Datasteru je v tomto prípade konzervatívne, pretože tu bolo pri webovom rozhraní overené, že nástroj pri rovnakom recepte, pri ktorom volanie knižnice zlyhá, reálne vygeneroval  čiastočný výstup.


Rovnako to platí aj pre kategóriu chýb, ako sebareferenčný FROM, ktorý v nástroji Dataster skončil pádom, ale príčina bola priradená samotnému programu, keďže táto chyba nebola zachytená vo validačnej a opravnej vrstve. 


Možné je aj spomenúť  model mistral:7b, kde kontrola padala hlavne v kategóriách na strane modelu (data\_too\_long: 168 a fk\_violation: 123). Tu je ukázané, že úspešnosť úrovne L3 zostáva nízka (29\%) aj po odrátaní zlyhaní pripísaných Datasteru. Metrika teda odlíši slabý model od obmedzení generátora.


Na záver je ale aj dôležité priznať limit pri priraďovaní chýb. Hraničné a generické scenáre (zero\_rows) sa nechávajú ako zlyhanie pipeline, čo znamená, že výsledná úspešnosť podhodnocuje pipeline.  Klasifikovanie chýb je taktiež úzko spojené s formátom chybových správ referenčného DBMS. Pri prechode na PostgreSQL bolo potrebné kategórie a priraďovanie chýb pozmeniť.

\section{Reprodukcia chýb}

Testovacie databázy \texttt{testdb21 - testdb23} boli vytvorené s úmyslom testovania správania generátora Dataster.

\begin{itemize}
    \item \textbf{\texttt{testdb21} Generovanie hodnôt pre finterval} Testovací scenár sa zameriaval na generovanie číselných rozsahov s plávajúcou rádovou čiarkou. Pri vygenerovaných SQL INSERT príkazoch sa objavili zlomky a tým sa splnil aj účel testovacieho scenára.
    \item \textbf{\texttt{testdb22} SMT jedinečnosť} Cieľom testovacieho scenára bola žiadosť o vytvorenie vyššieho počtu unikátnych hodnôt ako je SMT solver schopný vygenerovať čo vyvolá chybu smt\_count\_error.
    \item \textbf{\texttt{testdb23} generovanie hodnôt pre dinterval} výsledky ukázali, že pri generovaní hodnôt pre dátum sa vygenerované hodnoty približovali prevažne ku spodnej hranici dátumového intervalu.
\end{itemize}

\section{Demonštrácia}
Cieľom bolo overiť, či nástroj funguje aj mimo referenčného DBMS. Nástroj bol spustený nad 10 základnými databázami nad PostgreSQL v troch kolách pri teplote 0.0 s konfiguráciou mix (Tabuľka~\ref{tab:vysledky_demostracia}). Ukázalo sa, že 8 z 10 databáz boli plne funkčné s L3\_Adjusted 100\% na všetkých akceptovaných receptoch. 

Obe zlyhania (testdb4 - wrong FROM format, testdb7 - smt\_count) sú zhodné s MySQL, čiže problém nie je v prenositeľnosti, ale ide už o zdokumentované chyby nezávislé od DBMS. Pri klasifikovaní chýb bolo potrebné prispôsobiť klasifikátor formátu chybových správ PostgreSQL (odlišné chybové správy).

\label{sub:kap8_demonstracia}
\begin{table}[H]
    \centering
    \small
    \begin{tabular}{lc c l}
    \toprule
    \textbf{Databáza} & \textbf{L2} & \textbf{L3\_Adjusted} & \textbf{pozn.} \\
    \midrule
    testdb, 2, 5, 6, 8, 9, 10 & 3/3 & 100\% & zhodné s MySQL \\
    testdb3 & 3/3 & 100\% & 13$\times$ dataster\_finterval\_fraction \\
    testdb4 & -&N/A & Wrong FROM format (zhodné s MySQL) \\
    testdb7 & -&N/A & dataster\_smt\_count (zhodné s MySQL) \\
    \bottomrule
    \end{tabular}
    \caption{Demonštrácia na PostgreSQL (10 databáz, 3 kolá, teplota 0.0)}
    \label{tab:vysledky_demostracia}
\end{table}

\section{Zhrnutie}

V tejto kapitole boli vyhodnotené výsledky testov zamerané na otázky RQ1-RQ6 predstavené v Sekcii~\ref{sub:ciele}. V rámci RQ1 sa ukázalo, že teplota nemá dokázateľný vplyv na priemernú úspešnosť na databázy v testovacích sadách. Avšak pre teplotu 1.2 bola zaznamenaná najslabšia úspešnosť (na úrovni $L3_{Adjusted}$). Výsledky viedli k zvoleniu teploty 0.0 pre testovanie v ďalších fázach. 


Pri vyhodnotení RQ2 sa potvrdilo, že veľkosť modelu má výrazný vplyv na úspešnosť.  Tu konfigurácia 3b dosiahla najnižšiu priemernú úspešnosť a citlivosť na zložitosť schém spomedzi všetkých konfigurácií (7b, 3b a mix). Modely 7b a mix dosiahli pri L3\_Adjusted úspešnosť približujúcu sa k 100\%. Zaujímavý výsledok je schopnosť mix konfigurácie dosiahnuť podobný výsledok ako konfigurácia 7b, len pri nižších nárokoch.


Pre porovnanie rodín modelov v RQ3 platí, že llama3.1:8b dosiahla podobné výsledky ako referenčný qwen2.5:7b, zatiaľ čo mistral výrazne zaostával na oboch úrovniach. Tu sa ukázalo, že kvalitu úspechu nezaručuje len veľkosť modelu, ale rozhoduje aj rodina.


Pri RQ4 sa ukázalo, že každý sledovaný parameter (počet tabuliek, šírka tabuliek a hĺbka FK) sa správa inak a o konkrétnom zlyhaní rozhoduje charakter databázovej schémy. V rámci RQ5 boli identifikované body pri ktorých úspešnosť s rastúcou zložitosťou klesá. Úspešnosť klesá medzi 30 a 60 tabuľkami pri počte tabuliek, 100 a 200 stĺpcami pri šírke tabuliek a pri hĺbke FK reťazca sa zlom neukázal (pre konfiguráciu 7b a mix).


Na záver bolo pri otázke RQ6 potvrdené, že zlyhania programu možno spoľahlivo odlíšiť od defektov nástroja Dataster. Toto sa opiera o päť nezávislých pozorovaní.


Pre zhrnutie sa pri analýze výsledkov ukázalo, že lokálne LLM dokáže generovať prevažne vykonateľné použiteľné recepty, pričom hlavné limity úspešnosti ležia v troch oblastiach, a to zdokumentované defekty Dataster, identifikované medzery validačnej vrstvy pri sebareferenčnom odkaze FROM a kapacitné limity malého modelu na zložitých schémach.